%%
%% Test: Margin Notes — v1.1
%% Stage: wide margin
%% 
%% Purpose: Verify wide margin and margin notes:
%%   - \mnote{...}          — plain margin note
%%   - \marginfig{...}{...} — figure + caption in margin
%%   - Outer margin on odd/even pages (RTL)
%%   - Long text flow with many margin notes
%%   - Direction (RTL) on both odd and even pages
%% 

\documentclass{matinbook}

\title{تست حاشیه‌ی عریض — نسخه ۱.۱}
\author{تیم \lr{MatinBook}}
\date{۱۴۰۵}

\booktitle{تست حاشیه}

\begin{document}

\maketitle

\chapter{حاشیه‌ی عریض}
\label{chap:margin-test}

این سند، ویژگی حاشیه‌ی عریض \lr{\texttt{mb-layout}} را بررسی می‌کند.
هدف این است که ببینیم آیا یادداشت‌های حاشیه در صفحات زوج و فرد
به‌درستی جابجا می‌شوند، آیا متن یادداشت‌ها در هر دو حالت راست‌چین
است، و آیا ترکیب متن فارسی با یادداشت‌های حاشیه‌ای بدون مشکل
کار می‌کند.

%========================================
\section{یادداشت‌های صفحه‌ی اول}
\label{sec:page-1}
%========================================

این یک متن اصلی است که در آن از یادداشت حاشیه استفاده می‌کنیم.
\mnote{یادداشت ۱: این یک یادداشت کوتاه در صفحه‌ی اول است.}
و این ادامه‌ی متن اصلی است. همان‌طور که می‌بینید، متن اصلی
در ستون باریک‌تری قرار گرفته و فضای سمت بیرونی صفحه برای
یادداشت‌ها آزاد شده است.

در ادامه، یک یادداشت دیگر اضافه می‌کنیم.
\mnote{یادداشت ۲: این یادداشت کمی طولانی‌تر است تا ببینیم در چند خط شکسته می‌شود و آیا جهت متن درست است یا نه.}
و متن ادامه پیدا می‌کند. هدف این است که ببینیم یادداشت‌های
حاشیه چطور در کنار متن اصلی قرار می‌گیرند.

یادداشت‌های حاشیه معمولاً برای موارد زیر استفاده می‌شوند:
توضیح یک اصطلاح فنی، ارجاع به منبع، فرمول کمکی، هشدار،
یا نکته‌ی جانبی که در متن اصلی جای کافی ندارد.
\mnote{یادداشت ۳: در کتاب‌های ریاضی، فرمول‌های مهم اغلب در حاشیه تکرار می‌شوند.}
این سبک در کتاب‌های درسی آمریکایی مثل بویس و استوارت
بسیار رایج است.

متن اصلی ادامه دارد و ما یادداشت‌های بیشتری اضافه می‌کنیم.
\mnote{یادداشت ۴: یک یادداشت دیگر برای پر کردن صفحه.}
هدف این است که صفحه‌ی اول پر شود و به صفحه‌ی دوم برویم.

%========================================
\section{یادداشت‌های صفحه‌ی دوم (زوج)}
\label{sec:page-2}
%========================================

اکنون در صفحه‌ی دوم هستیم. این صفحه باید زوج باشد و
یادداشت‌ها باید در حاشیه‌ی مقابل صفحه‌ی اول قرار بگیرند.

\mnote{یادداشت ۵: یادداشت در صفحه‌ی دوم (زوج).}
متن اصلی اینجا ادامه دارد. باید ببینیم که حاشیه در صفحه‌ی
بعد به سمت مقابل منتقل می‌شود یا نه.

\mnote{یادداشت ۶: یک یادداشت طولانی‌تر در صفحه‌ی زوج. این یادداشت باید در چند خط شکسته شود و جهت متن آن باید راست‌چین باشد.}
متن بیشتر... متن بیشتر... متن بیشتر... متن بیشتر...
متن بیشتر... متن بیشتر... متن بیشتر... متن بیشتر...
متن بیشتر... متن بیشتر... متن بیشتر... متن بیشتر...

\mnote{یادداشت ۷: یادداشت کوتاه در صفحه‌ی زوج.}
متن اصلی ادامه دارد.

\mnote{یادداشت ۸: این یادداشت دارای \lr{\textbf{bold}} و \lr{\emph{italic}} است.}
و یک یادداشت دیگر با فرمول ریاضی:
\mnote{یادداشت ۹: فرمول $E = mc^2$ در یادداشت.}
متن اصلی ادامه دارد.

%========================================
\section{تصویر در حاشیه}
\label{sec:marginfig}
%========================================

اکنون یک تصویر در حاشیه قرار می‌دهیم.
\marginfig{example-image}{نمودار نمونه‌ی حاشیه}
و متن اصلی در ستون اصلی ادامه پیدا می‌کند.

%========================================
\section{ترکیب با محیط‌های علمی}
\label{sec:with-theorems}
%========================================

حاشیه‌ی عریض باید با محیط‌های علمی هم سازگار باشد.

\begin{theorem}[قضیه فیثاغورس]
    در هر مثلث قائم‌الزاویه، مربع وتر برابر است با مجموع
    مربعات دو ضلع دیگر:
    \[
    a^2 + b^2 = c^2
    \]
    \label{thm:pythagoras}
\end{theorem}

\mnote{یادداشت ۱۰: این قضیه یکی از قدیمی‌ترین قضایای ریاضی است.}

\begin{proof}
    اثبات این قضیه با روش‌های مختلفی امکان‌پذیر است.
    \label{prf:pythagoras}
\end{proof}

\mnote{یادداشت ۱۱: بیش از ۳۷۰ اثبات برای این قضیه ثبت شده است.}

\begin{example}[حل معادله درجه دوم]
    معادله $x^2 - 5x + 6 = 0$ را حل کنید.
    \label{ex:quadratic}
\end{example}

\mnote{یادداشت ۱۲: برای حل این معادله می‌توان از تجزیه یا فرمول عمومی استفاده کرد.}

\begin{solution}
    با تجزیه داریم:
    \[
    x^2 - 5x + 6 = (x - 2)(x - 3) = 0
    \]
    پس $x = 2$ یا $x = 3$.
\end{solution}

\mnote{یادداشت ۱۳: فرمول عمومی $x = \frac{-b \pm \sqrt{b^2 - 4ac}}{2a}$ هم جواب یکسان می‌دهد.}

%========================================
\section{یادداشت‌های پشت سر هم}
\label{sec:consecutive-notes}
%========================================

یادداشت‌های حاشیه می‌توانند پشت سر هم بیایند.

\mnote{یادداشت ۱۴: یادداشت اول.}
متن اصلی. \mnote{یادداشت ۱۵: یادداشت دوم.} متن اصلی.
\mnote{یادداشت ۱۶: یادداشت سوم.} متن اصلی ادامه دارد.
\mnote{یادداشت ۱۷: یادداشت چهارم.}
متن اصلی. \mnote{یادداشت ۱۸: یادداشت پنجم.}

%========================================
\section{یادداشت‌های طولانی}
\label{sec:long-notes}
%========================================

\mnote{یادداشت ۱۹: این یک یادداشت نسبتاً طولانی است که در چند خط نوشته شده تا ببینیم آیا در حاشیه به‌درستی شکسته می‌شود و آیا جهت متن در همه‌ی خطوط راست‌چین می‌ماند یا نه.}
متن اصلی ادامه دارد. اگر یادداشت خیلی بلند باشد، ممکن است
به صفحه‌ی بعد سرریز کند.

\mnote{یادداشت ۲۰: یادداشت پایانی.}
متن اصلی. این آخرین یادداشت این بخش است.

%========================================
\section{کد و اصطلاحات لاتین}
\label{sec:latin}
%========================================

در این بخش، چند کلمه و دستور لاتین را داخل متن فارسی
می‌آوریم تا مطمئن شویم \lr{\texttt{\textbackslash lr}} درست کار می‌کند.

دستور \lr{\texttt{\textbackslash mnote}} برای یادداشت حاشیه استفاده می‌شود.

دستور \lr{\texttt{\textbackslash marginfig}} برای تصویر در حاشیه است.

پکیج \lr{\texttt{sidenotesplus}} برای مدیریت یادداشت‌ها به کار می‌رود.

پارامتر \lr{\texttt{marginparwidth}} عرض ناحیه‌ی یادداشت را تعیین می‌کند.

پارامتر \lr{\texttt{marginparsep}} فاصله‌ی بین متن و یادداشت را تعیین می‌کند.

%========================================
\section{نتیجه‌گیری}
\label{sec:conclusion}
%========================================

اگر این سند بدون خطا کامپایل شود، یعنی ویژگی‌های زیر
به‌درستی کار می‌کنند:

\lr{\texttt{\textbackslash mnote}} یادداشت حاشیه می‌سازد،
\lr{\texttt{\textbackslash marginfig}} تصویر با زیرنویس در حاشیه
می‌سازد، حاشیه‌ی عریض در صفحات زوج و فرد درست جابجا می‌شود،
متن اصلی باریک‌تر شده و فضای بیشتری برای یادداشت‌ها دارد،
و ترکیب یادداشت حاشیه با محیط‌های علمی مثل قضیه، اثبات،
مثال، و راه‌حل بدون مشکل کار می‌کند.

\textbf{وضعیت تست:} قبول.

\end{document}